% 正文按完整句子换行。参考文献、图注、作者简介不计入正文篇幅。
一家连锁商店想比较各门店的净销售额，看起来只需做一次加减法。
但订单、退款和门店信息往往分散在不同表格中，日期写法不一，同一家门店可能使用不同简称，一笔订单还可能对应多笔退款。
假设一笔100元的订单分两次退款，金额分别为20元和10元。
如果把订单与退款记录直接关联，再对每行相减求和，就可能把100元算两遍，得到170元，而不是正确的70元。
程序没有报错，图表也很漂亮，结论却已经偏离事实。

分析之前，必须先弄清数据的含义、组织方式和使用规则。
找到相关数据，识别问题，连接不同来源，完成必要转换，再检查结果，这些工作共同构成了数据准备。
它不是分析过程中的附属杂务，而是把原始记录变成可靠依据的关键环节。
智能数据准备希望进一步降低这项工作的门槛：让人说明目标，让系统协助理解数据、设计流程、执行检查，并在有疑问时寻求确认，而不是要求使用者从头编写每一行程序。\cite{demystify}

\section{数据准备，准备的究竟是什么}

数据清洗只是数据准备的一部分。
清洗主要处理重复、缺失、错误和冲突；准备还包括发现数据、解释字段、整合来源、改变组织形式，以及为学习任务补充必要的标注。
例如，将多个月份分别占一列的报表改成按月份逐行记录，是结构转换；判断两个不同名称是否指向同一家企业，是实体匹配；把订单与退款组织到同一个统计口径下，则需要同时考虑关联与计算。
这些操作往往相互影响，不能机械地按固定顺序执行。\cite{demystify,quality}

准备对象也不局限于电子表格。
合同里的条款、图片中的文字、设备日志中的时间记录，常需要先提取并组织，才能与数据库中的信息共同使用。
提取过程本身也可能出错，因此既要保存整理后的结果，也要保留能够返回原文、原图或原始记录的位置。
这使数据准备同时承担了连接不同信息形式和保留证据的任务。\cite{demystify,datasheets}

准备也不是分析前只做一次的工作。
当源系统增加字段、门店改变名称，或统计规则发生调整，原有流程可能不再适用。
可持续的数据准备需要识别这些变化，重新检查受影响的步骤，并保留不同版本之间的关系。
否则，一次正确的结果并不能保证下个月仍然正确。\cite{quality}

“准备好”也没有脱离用途的统一答案。
分析门店经营，需要保留退货信息；识别异常交易，恰恰不能先删掉看起来不正常的订单。
同一个空白值，可能表示尚未采集、不适用，也可能意味着系统漏记，不能一律填成零。
数据质量因此既包括准确性、完整性和一致性，也取决于它是否及时、是否适合当前任务。
把表格整理得整齐，并不等于把事实解释正确。\cite{quality}

智能之处首先在于理解这种“用途与处理方式”的关系。
传统程序按明确规则执行，语言模型则能够利用列名、文本内容和任务描述判断哪些操作可能需要进行。
但语言模型的判断并不是事实本身。
比较可靠的分工是：模型负责理解与提出方案，数据库和经过检验的程序负责运算，质量检查负责发现偏差，人负责确认业务含义和重要决策。
因此，给工具加上聊天窗口只是入口，能否根据真实执行结果修正方案，才是更重要的能力。\cite{fm,autoprep,deepprep}

\articlefigure{fig01_scope}{从原始数据到可信结果。数据准备围绕具体用途展开，发现、理解、清洗、整合、转换与验证可以反复交织，来源和变更记录贯穿全过程。}{fig:scope}

\section{它与我们的生活和工作有什么关系}

\textbf{在办公与经营中，减少重复劳动只是第一步。}
合并月报、核对库存、分析客户反馈，都需要处理来源不同、结构不齐的记录。
智能系统可以建议日期转换、字段对应和统计步骤，把一次处理保存为可重用的流程。
真正重要的收益，是让使用者有机会看清计算依据，而不是仅仅更快得到一个数字。
开篇的退款例子中，系统应先确认统计规则，再判断是否需要按订单汇总退款，最后检查销售额是否被重复计算。
这也是自然语言驱动的数据准备研究试图解决的实际问题。\cite{autoprep,deepprep}

\textbf{在科研和公共服务中，准备数据要保留解释所需的背景。}
把不同设备的观测记录放在一起，不能只统一文件格式，还要确认单位、采样时间和测量条件。
罕见的高值可能是传感器故障，也可能是真实的极端事件，删除之前需要查证。
医疗研究中的不同编码和记录方式同样需要谨慎对应，自动处理不能替代专业判断。
科学数据管理强调可查找、可访问、可互操作和可重用，背后的要求正是让后来使用数据的人理解它的来源、含义与适用范围。\cite{fair}

\textbf{在大模型训练中，数据准备影响模型学到什么。}
网页、文档、图片和问答材料需要去重、筛选、整理来源，并检查不同内容的覆盖情况。
只追求数量，可能保留大量重复信息；过度追求“干净”，又可能删去少见但有价值的表达。
Data-Juicer等系统把处理步骤组织成可组合的操作，并连接数据分析和模型评估，使研究者能够比较不同处理方案，而不只是凭经验选择材料。\cite{datajuicer,llmdata}
这里有两个不同方向：用人工智能帮助准备数据，以及为人工智能准备数据。
两者相互促进，但不能混为一谈。

\textbf{在机器人学习中，还要保证不同记录描述的是同一时刻、同一个动作。}
相机画面、机械臂位置和控制指令即使各自正确，时间没有对齐，也可能形成错误的学习样本。
不同机器人还可能使用不同的动作表示。
Open X-Embodiment对多种机器人数据进行汇集和标准化，说明跨设备学习离不开数据组织工作。\cite{oxe}
让智能系统进一步协助检查时间同步、动作含义和样本覆盖，是值得探索的方向，而不是已经普遍解决的问题。
图\ref{fig:applications}概括了这些用途，它们的共同点是：数据准备必须服务于具体任务，而非套用同一套清洗规则。

\articlefigure{fig02_applications}{数据准备服务于不同用途。办公经营关注口径和关联，科研关注来源与测量条件，大模型训练关注质量与覆盖，机器人学习还需要时间和动作的一致性。图中为用途示意，并非特定系统的部署案例。}{fig:applications}

\section{研究正在从单项工具走向完整流程}

\subsection{让机器帮助人整理数据}

数据准备并不是大模型出现后才开始的研究。
早期系统主要依靠规则与脚本，随后出现了更易操作的交互式工具。
Wrangler让用户通过表格界面示范或选择变换，再生成相应处理步骤，降低了逐行编程的负担。\cite{wrangler}
HoloClean进一步综合数据约束、统计关系和参考信息推断可能的修复结果，说明机器能够帮助发现仅靠手写规则难以覆盖的问题。\cite{holoclean}

另一条路线是自动寻找转换程序。
Auto-Tables针对结构不规则的表格合成多步转换，并不要求用户先给出转换后的示例表。\cite{autotables}
这类研究提醒我们，自动化不只有“大模型写代码”一种实现方式。
规则、程序搜索和统计学习各有所长，至今仍是智能系统的重要组成部分。

语言模型带来的新变化，是能够较自然地理解字段含义和任务描述。
早期探索已经将实体匹配、错误检测等数据任务转化为语言问答；Jellyfish则通过针对性训练，使可本地运行的模型支持多类数据预处理任务。\cite{fm,jellyfish}
不过，做好一次匹配或一次修复，不等于能稳定完成一条相互依赖的长流程。
研究的重心因而开始从单项能力转向多步骤协同。

\subsection{从理解问题，到根据执行结果回退修正}

AutoPrep关注一个容易被忽视的事实：同一张表面对不同问题，需要的准备工作并不相同。
它把任务分为规划、编程和执行，让系统先根据问题确定需要哪些高层操作，再转成具体程序。
例如，回答某年的年龄分布可能需要从出生日期计算年龄，而回答地区分布则需要识别和规范化地域信息。
这项工作的价值，是把数据准备与实际问题连接起来，而非不加区分地清洗所有字段。\cite{autoprep}

DeepPrep进一步处理多表、多步骤准备中“错误很晚才暴露”的困难。
它保留中间表和操作反馈，通过树状的过程组织方式保存不同尝试，使系统可以退回较早的决策，而不只是反复修改最后一段代码。
它还通过分阶段训练，让模型学习操作使用、执行反馈和回退行为。
在论文所测任务中，这种专门训练改善了开源模型的准确率与推理成本权衡，但这不意味着任意业务数据都能无人值守地处理。\cite{deepprep}

图\ref{fig:agent}用退款例子说明这种区别。
当系统发现销售额重复，真正需要修正的可能是前面的关联方式，而不是后面的加减公式。
能够保存现场、找到较早的错误步骤并重新尝试，是数据准备智能体区别于一次性代码生成的重要方向。
这里的“智能体”，指能够围绕目标调用工具、观察结果并继续行动的系统，不意味着它已经拥有人的全部判断能力。

\articlefigure{fig03_agent}{执行反馈与回退修正。示例中，一笔100元订单对应20元和10元两笔退款。直接关联后求和可能得到错误的170元；先汇总退款、再关联计算，才能得到70元。数字为解释原理而设，不是实验结果。}{fig:agent}

\subsection{让方法适应新数据，也让能力能够被检验}

另一项挑战是换一个行业、换一批表格后，系统能否继续工作。
KnowTrans通过复用已有任务的知识，并引入外部信息，帮助专用数据准备模型以较少样本适应新数据和新任务。
它并非证明模型已经无需领域知识，恰恰说明恰当的知识补充仍然重要。\cite{knowtrans}

Text-to-Pipeline把自然语言要求转化为处理流程作为明确任务，并构建PARROT基准，包含约1.8万项任务和16类核心操作。
其研究发现，模型不仅会弄错操作之间的组合关系，也会把文字要求对应到错误的列、条件或参数。
配套的Pipeline-Agent利用中间执行状态改进流程，但仍存在明显错误。\cite{texttopipeline}
BAT则探索在没有目标表实例的条件下，通过受控操作和树搜索合成准备流程。\cite{bat}
这些工作与前述研究从不同入口推进同一问题：如何把大致合理的意图，落实成可以检查的操作。

总体而言，现阶段既不是“数据准备已经全面自动化”，也不是“模型只能做演示”。
单项任务已有丰富方法，较明确的多步骤任务正在取得进展，跨领域、含糊需求和长期运行仍是薄弱环节。
PrepBench将意图澄清、代码生成和工作流表达纳入评测，也发现现有系统仍难稳定贯通这些能力。\cite{prepbench}
不同论文采用的数据、任务和成本口径并不相同，不能把各自的分数直接排成一张通用排行榜。

\section{从能运行到可信赖，还需要跨过哪些门槛}

\textbf{首先是说清目标，而不是让模型猜口径。}
“统计活跃用户”究竟按登录、购买还是互动计算？
“删除重复”是删除完全相同的行，还是合并同一对象的多条记录？
这些歧义会改变结果。
在动手转换之前，还要确认选中的数据是否来自正确的时间段，字段说明是否对应当前版本。
数据目录、业务词典和经过确认的历史流程可以提供帮助，但系统应展示采用了哪份说明，避免把过时经验直接套到新任务上。
系统应该优先追问少数影响最大的条件，并在执行前展示拟采用的定义。
PrepBench的结果表明，交互澄清能够帮助改进表现，但问题问得多不等于问得好。\cite{prepbench}

\textbf{其次是把纠错建立在可观察的证据上。}
程序成功运行，只能说明它完成了计算，不能说明计算符合原意。
关联前后记录数为何变化，关键标识是否仍然唯一，某类样本是否被大量过滤，都应进入检查过程。
中间状态和变更记录可以帮助追溯问题，而独立的规则、统计检查与人工抽查应与模型判断相互补充。
执行反馈和回退机制为此提供了基础，但验证标准本身仍需要领域知识。\cite{quality,deepprep}

\textbf{第三是承认无法恢复的信息。}
模型可以根据上下文提出缺失值的候选答案，却不能保证猜测等于真实记录。
对难以确认的内容，应保留缺失标记或多个候选，并说明依据，而不是把不确定性藏进一张看似完整的表里。
处理异常值也应区分错误与少见事件，避免“修好”数据的同时抹掉研究对象。
数据集说明文档所强调的来源、组成、处理过程和使用限制，有助于把这些边界传递给后续使用者。\cite{datasheets}

\textbf{第四是让可靠性与规模、成本一起设计。}
面对海量记录，不宜让语言模型逐行处理所有内容。
可以先用数据库完成计数、分组和确定性的格式转换，再把难以解释的字段、少量代表样本和检查摘要交给模型。
已确认的步骤可以缓存和重用，数据变化后只重新处理受影响的部分。
但抽样可能漏掉稀有错误，因此必须配合覆盖全量数据的规则检查和有针对性的复核。
这种分工比单纯扩大模型更符合数据系统的特点。\cite{demystify,datajuicer}

\textbf{第五是把权限和责任放进流程。}
系统不应因为“能够调用工具”，就默认可以读取全部记录、向外部服务发送敏感内容或覆盖原始数据。
本地运行可以减少数据外传的路径，却不能替代访问控制和操作审计。\cite{jellyfish}
合理的设计是保留只读原始副本，限制执行权限，记录关键修改，并把不可逆操作交给人确认。
交付物也不应只有结果表，还应包括处理流程、验证依据和适用范围。

\textbf{最后是评价真正的使用效果。}
好的准备流程不仅要生成正确结果，还要控制时间、计算开销和人工修订量。
面向模型训练，还应检查不同数据群体的覆盖和下游表现；面向经营分析，则要检查口径变化是否影响结论。
评测时必须避免把测试答案或未来信息混入准备过程，也不能用总体分数掩盖少数高风险失败。
基准测试有助于比较方法，真实应用仍需要结合自己的数据、风险和长期维护成本做验证。\cite{prepbench,llmdata}
图\ref{fig:trust}将这些要求归纳为相互配合的检查环节。

\articlefigure{fig04_trust}{可信数据准备的关键检查。目标、执行、事实、成本、权限与效果需要共同约束系统行为；自动化程度应随风险调整，而不是把所有任务一律交给模型。}{fig:trust}

\section{未来：从一次性整理走向持续的数据协作}

更可行的发展路径，是逐步扩大经过验证的自动化范围。
低风险、规则明确、重复发生的工作，可以更多交给系统执行；涉及口径变化、敏感信息和重要决策的步骤，则保留人的确认。
这并非自动化的退让，而是把人的专业判断用在最值得介入的地方。
能够说明何时需要帮助，与能够独立完成任务同样重要。

另一个方向，是让准备过程与后续分析相互反馈。
当分析发现某类记录持续缺失，或模型在某些场景中频繁失败，系统可以重新检查数据覆盖、标注方式和处理规则。
DeepAnalyze等数据科学智能体已尝试连接数据处理、分析与报告生成，为这种衔接提供了探索。\cite{deepanalyze}
但更长的流程也意味着错误可能传播得更远，因此每个阶段仍需要独立检查，不能让一份表达流畅的报告替代对数据的核实。

随着多模态和跨设备应用发展，数据准备还需要更好地描述时间、空间、动作和来源之间的关系，并在数据持续更新时保持一致。
可复用的操作定义、数据说明和公开评测，有望减少重复建设，使不同系统的结果更容易交换和核验。\cite{fair,oxe,datasheets}
这些是需要继续推进的研究方向，而不是对短期全面无人化的承诺。

智能数据准备最终要解决的，不只是“让机器替人整理数据”，而是让数据的用途更明确、处理更可靠、错误更容易被发现。
当系统能够同时交付可用的数据、可重现的过程和可检查的依据，人工智能才获得了更值得信任的数据基础。
